\documentclass[11pt, twocolumn]{article}
\usepackage[margin=0.65in, columnsep=0.25in, bottom=0.45in]{geometry}
\usepackage{newtxtext}
\usepackage{titlesec}
\usepackage{enumitem}
\usepackage{xurl}
\usepackage{hyperref}
\usepackage{xcolor}

% Tight compacting
\setlist{nosep, leftmargin=*} 
\linespread{0.93} 
\setlength{\parskip}{2.5pt} 
\renewcommand{\footnotesize}{\tiny}

% Tight section formatting with visual distinction
\titleformat{\section}{\large\bfseries\scshape}{\thesection}{1em}{}[\vspace{1pt}\titlerule]
\titlespacing*{\section}{0pt}{6pt}{3pt}

\hypersetup{
    colorlinks=true,
    linkcolor=black,
    urlcolor=blue,
    citecolor=black,
    breaklinks=true,
}

\pagestyle{plain} 

% Compact Bibliography
\makeatletter
\renewenvironment{thebibliography}[1]
     {\section*{References}%
      \tiny
      \setlength{\topsep}{0pt}
      \list{\@biblabel{\@arabic\c@enumiv}}%
           {\settowidth\labelwidth{\@biblabel{#1}}%
            \leftmargin\labelwidth
            \advance\leftmargin\labelsep
            \itemsep0pt\parsep0pt\parskip0pt
            \usecounter{enumiv}%
            \let\p@enumiv\@empty
            \renewcommand\theenumiv{\@arabic\c@enumiv}}%
      \sloppy
      \clubpenalty4000
      \@clubpenalty \clubpenalty
      \widowpenalty4000%
      \sfcode`\.\@m}
     {\def\@noitemerr
       {\@latex@warning{Empty `thebibliography' environment}}%
      \endlist}
\makeatother

\begin{document}

\twocolumn[
  \begin{@twocolumnfalse}
    \noindent
    \textbf{\LARGE AI Reading Exercise: Class 7} \hfill \textbf{Daniel Hardesty Lewis} \\
    \noindent
    \textit{AI, Energy, and Climate} \hfill March 3, 2026 \\
    \textbf{PLAN A6613: AI and the Future of Cities} \hfill Spring 2026
    \vspace{3pt}
    \hrule
    \vspace{6pt}
  \end{@twocolumnfalse}
]

\section*{Tool \& Process}
I used \textbf{Claude Opus 4.6} with extended thinking via the Antigravity IDE.\footnote{Verbatim prompt log, downloaded sources, annotations, and a claim-by-claim verification audit are archived at \url{https://github.com/dhardestylewis/plan-a6613-ai-reading-class-3/tree/main/week7}.} I adopted a \textit{source-first} methodology, downloading Google's 2024 and 2025 Environmental Reports and the IEA's data center energy analysis before prompting the LLM, so that I could compare its claims against primary documents immediately. This approach proved essential because Google's sustainability reporting uses an accounting methodology that can only be evaluated by cross-referencing corporate disclosures against independent analyses. The prompt asked Claude to answer the assignment question verbatim. Claude assembled accurate emission figures from both reports but defaulted to Google's preferred \textit{market-based} accounting, which credits renewable energy purchases against grid emissions. It did not distinguish this from \textit{location-based} accounting, which reflects the actual carbon intensity of the power grids where Google's data centers operate. Surfacing that distinction, and tracing the conditional framing of Google's efficiency claims to their source papers, required manual cross-referencing across all downloaded sources.

\section*{Key Findings}
\noindent Google's efficiency evidence is self-authored, pre-GenAI, and conditional. Google discloses this, but AI tools strip the provenance when reproducing these claims, presenting them as independently verified current achievements. Detecting this requires the manual cross-referencing that AI tools cannot perform.

\textbf{\textit{1. Emissions have risen sharply despite efficiency gains.}}
Google's 2024 Environmental Report disclosed 14.3 million metric tons of CO2 equivalent in 2023, a 13\% year-over-year increase and 48\% above its 2019 baseline (Google, 2024). The 2025 report, covering fiscal year 2024, shows a further rise to a cumulative 51\% increase since 2019, driven by a 27\% jump in data center electricity consumption in a single year (Google, 2025). Google frames this around market-based accounting, which incorporates renewable energy credits (Google, 2024; Google, 2025). An independent analysis using \textit{location-based} accounting, which measures actual grid emissions at the point of consumption, found a 65\% increase over the same period (Milman, 2025). The gap between these two figures reveals how renewable energy purchases can mask the physical carbon intensity of a company's operations.

\textbf{\textit{2. Efficiency claims are real but self-sourced and overwhelmed by growth.}}
Google reports that its sixth-generation TPU, Trillium, is 67\% more energy-efficient than its predecessor, and that combining best practices can reduce training energy by 100 times and emissions by 1,000 times (Google, 2024). Google labels these as ``our research,'' and the underlying paper confirms this: its authors are all Google employees, it predates the generative AI boom by seven months, and its claims are explicitly conditional on applying all practices simultaneously (Patterson et al., 2022). Google's data centers achieve a power usage effectiveness of 1.10, compared with the industry average of 1.58 (Uptime Institute, 2023), and the median energy per Gemini text prompt fell 33-fold in twelve months (Google, 2025). Yet the IEA documents that combined electricity use by Amazon, Microsoft, Google, and Meta more than doubled between 2017 and 2021, reaching approximately 72 terawatt-hours, and continues to grow 20 to 40 percent annually (IEA, 2023). Per-unit efficiency gains, however real, are not keeping pace with absolute growth.

\section*{Verification and Reflection}
Claude correctly identified both Environmental Reports and cited emission figures accurately. Its limitation was \textit{framing}, not fabrication: it stripped the provenance that Google itself provides. Where Google writes ``our research,'' Claude presented the efficiency claims without noting the self-sourcing, the pre-GenAI timing, or the conditionality. It defaulted to market-based accounting without surfacing the location-based alternative. Google's emissions have risen 51\% since 2019 by its own market-based measure and 65\% by location-based accounting, yet the IEA's net-zero scenario requires data center emissions to \textit{halve} by 2030 (IEA, 2023). The source-first methodology recovered the provenance and surfaced this contradiction; the AI tool, left to its own framing, did not.

\begin{thebibliography}{9}
\bibitem{google2024} Google. (2024, July 2). \textit{Google's 2024 Environmental Report}. \url{https://blog.google/outreach-initiatives/sustainability/2024-environmental-report/}
\bibitem{google2025} Google. (2025). \textit{Google's 2025 Environmental Report}. \url{https://www.gstatic.com/gumdrop/sustainability/google-2025-environmental-report.pdf}
\bibitem{iea} International Energy Agency. (2023). Data centres and data transmission networks. \url{https://www.iea.org/energy-system/buildings/data-centres-and-data-transmission-networks}
\bibitem{milman} Milman, O. (2025, February 12). Google's carbon emissions have risen 65\% since 2019, analysis finds. \textit{The Guardian}. \url{https://www.theguardian.com/technology/2025/feb/12/google-carbon-emissions-rise-ai}
\bibitem{patterson} Patterson, D., Gonzalez, J., Hölzle, U., et al. (2022). The carbon footprint of machine learning training will plateau, then shrink. \textit{Computer}, 55(7), 18--28.
\bibitem{uptime} Uptime Institute. (2023). \textit{2023 Global Data Center Survey}. \url{https://uptimeinstitute.com/resources/research-and-reports/uptime-institute-global-data-center-survey-results-2023}
\end{thebibliography}

\end{document}
